\documentclass[11pt]{article}

\usepackage[french, english]{babel}

\usepackage[T1]{fontenc}

\usepackage{amsfonts}
\usepackage{amsmath}
\usepackage{amssymb}

\usepackage{graphicx}
\usepackage{multirow}
\usepackage{multicol}
\usepackage{xcolor}
\usepackage[shortlabels]{enumitem}

\usepackage[margin=2cm]{geometry}

\usepackage{mathtools}
\usepackage{siunitx}
\usepackage{braket}

\usepackage{listings}

\usepackage{tikz}
\usetikzlibrary{shapes.geometric, arrows.meta, positioning}

\tikzstyle{block} = [rectangle, rounded corners, minimum width=4cm, minimum height=1cm,
text centered, draw=black, fill=blue!10]
\tikzstyle{process} = [rectangle, rounded corners, minimum width=5cm, minimum height=1cm,
text centered, draw=black, fill=green!10]
\tikzstyle{io} = [trapezium, trapezium left angle=70, trapezium right angle=110,
minimum width=4cm, minimum height=1cm, text centered, draw=black, fill=orange!15]
\tikzstyle{decision} = [diamond, aspect=2, draw=black, fill=yellow!15, text centered]
\tikzstyle{arrow} = [thick, ->, >=Stealth]

\renewcommand\lstlistingname{Code Python}

\lstset{frame=tb,
  language=Python,
  numbers=left,
  stepnumber=1,
  aboveskip=3mm,
  belowskip=3mm,
  showstringspaces=false,
  columns=flexible,
  basicstyle={\small\ttfamily},
%   numbers=none,
  numberstyle=\tiny\color{gray},
  keywordstyle=\color{blue},
  commentstyle=\color{paired2b},
  stringstyle=\color{paired3b},
  breaklines=true,
  breakatwhitespace=true,
  tabsize=1,
  frame=single,
%   showtabs=false,
}

\newcounter{questioncounter}
\newenvironment{question}[1]{
\refstepcounter{questioncounter}
\noindent\rule{0.05\textwidth}{1pt}\,\,\raisebox{-.3\ht\strutbox}{\textbf{Question \thequestioncounter : #1}}\,\,\hrulefill\par\vspace{0.4cm}
}{\par}

\begin{document}

%\renewcommand*{\contentsname}{Table des matières}

\hrule

\begin{center}
  {\Large PHQ334 - Mécanique quantique} \\ \vspace{5mm}
  {\LARGE\sffamily Devoir 5} \\
  Makram Abid, Alexis Jean\\
  À remettre le 8 décembre 2025
\end{center}

\hrule

%?\tableofcontents

\vspace{1cm}

\begin{question}{Interaction Heisenberg isotrope}
\subsection*{a) Matrices des opérateurs ($S_\mu^1$) (avec ($\mu=z,+,-,x,y$))}
Puisque nous savons que la base $\{{\ket{\uparrow\uparrow},\ \ket{\uparrow\downarrow},\ \ket{\downarrow\uparrow},\ \ket{\downarrow\downarrow}}\}$ nous donne la matrice identité 4x4, on peut écrit ($S_\mu^1=S^\mu\otimes\mathbb I$) et ($S_\mu^2=\mathbb I\otimes S^\mu$).
Dans la base des spins donnée plus haut,
\[
S_z^1=\frac{\hbar}{2}\operatorname{diag}(1,1,-1,-1),
\qquad
S_z^1=\frac{\hbar}{2}\operatorname{diag}(1,-1,1,-1).
\]
Pour les opérateurs +/-, rappel : $(S_+=\hbar\begin{pmatrix}0&1\\0&0\end{pmatrix})$, $(S_-=\hbar\begin{pmatrix}0&0\\1&0\end{pmatrix})$. Ainsi
\[
S_+^1=\hbar\begin{pmatrix}
0&1&0&0\\
0&0&0&0\\
0&0&0&1\\
0&0&0&0\\
\end{pmatrix},\qquad
S_-^1=\hbar\begin{pmatrix}
0&0&0&0\\
1&0&0&0\\
0&0&0&0\\
0&0&1&0\\
\end{pmatrix},
\]
et pour le spin 2 (opérant sur la seconde composante)
\[
S_+^2=\hbar\begin{pmatrix}
0&0&1&0\\
0&0&0&1\\
0&0&0&0\\
0&0&0&0\\
\end{pmatrix},\qquad
S_-^2=\hbar\begin{pmatrix}
0&0&0&0\\
0&0&0&0\\
1&0&0&0\\
0&1&0&0\\
\end{pmatrix}.
\]
Enfin, $(S_x=\tfrac12(S^+ + S^-))$, $(S_y=\tfrac{1}{2i}(S^+ - S^-))$, donc on peut écrire explicitement $(S_x^1,S_y^1,S_x^2,S_y^2)$ en faisant les opérations appropriés avec les matrices trouvées plus haut t.q. :\\
Pour $S_x^1$ :
\[
S_x^1=\frac{\hbar}{2}\begin{pmatrix}
0&1&0&0\\
0&0&0&0\\
0&0&0&1\\
0&0&0&0\\
\end{pmatrix}
+ \begin{pmatrix}
0&0&0&0\\
1&0&0&0\\
0&0&0&0\\
0&0&1&0\\
\end{pmatrix}
= \frac{\hbar}{2}\begin{pmatrix}
0&1&0&0\\
1&0&0&0\\
0&0&0&1\\
0&0&1&0\\
\end{pmatrix}
\]
Pour $S_y^1$ :
\[
S_y^1=\frac{\hbar}{2i}\begin{pmatrix}
0&1&0&0\\
0&0&0&0\\
0&0&0&1\\
0&0&0&0\\
\end{pmatrix}
- \begin{pmatrix}
0&0&0&0\\
1&0&0&0\\
0&0&0&0\\
0&0&1&0\\
\end{pmatrix}
= \frac{\hbar}{2i}\begin{pmatrix}
0&1&0&0\\
-1&0&0&0\\
0&0&0&1\\
0&0&-1&0\\
\end{pmatrix}
\]
Pour $S_x^2$ :
\[
S_x^2=\frac{\hbar}{2}\begin{pmatrix}
0&0&1&0\\
0&0&0&1\\
0&0&0&0\\
0&0&0&0\\
\end{pmatrix}
+ \begin{pmatrix}
0&0&0&0\\
0&0&0&0\\
1&0&0&0\\
0&1&0&0\\
\end{pmatrix}
= \frac{\hbar}{2}\begin{pmatrix}
0&0&1&0\\
0&0&0&1\\
1&0&0&0\\
0&1&0&0\\
\end{pmatrix}
\]
Pour $S_y^2$ :
\[
S_y^2 = \frac{\hbar}{2i}\begin{pmatrix}
0&0&1&0\\
0&0&0&1\\
0&0&0&0\\
0&0&0&0\\
\end{pmatrix}
- \begin{pmatrix}
0&0&0&0\\
0&0&0&0\\
1&0&0&0\\
0&1&0&0\\
\end{pmatrix}
= \frac{\hbar}{2i}\begin{pmatrix}
0&0&1&0\\
0&0&0&1\\
-1&0&0&0\\
0&-1&0&0\\
\end{pmatrix}
\]

\subsection*{(b) Représentation de l'hamiltonien}
On peut décomposer les matrices S selon la forme :
\[
H_H = J\left(
S_x^1S_x^2 + S_y^1 S_y^2 + S_z^1 S_z^2
\right).
\]
On calcule séparément chaque terme t.q. :
\[
S_z^1 S_z^2 =
\frac{\hbar^2}{4}
\begin{pmatrix}
1 & 0 & 0 & 0\\
0 & -1 & 0 & 0\\
0 & 0 & -1 & 0\\
0 & 0 & 0 & 1\\
\end{pmatrix}.
\]
\[
S_+^1 S_-^2 =
\hbar^2
\begin{pmatrix}
0 & 0 & 0 & 0\\
0 & 0 & 0 & 0\\
0 & 1 & 0 & 0\\
0 & 0 & 0 & 0\\
\end{pmatrix},
\quad
S_-^1 S_+^2 =
\hbar^2
\begin{pmatrix}
0 & 0 & 0 & 0\\
0 & 0 & 1 & 0\\
0 & 0 & 0 & 0\\
0 & 0 & 0 & 0\\
\end{pmatrix}.
\]
On sait que :
\[
S_x^1 S_x^2 + S_y^1 S_y^2
= \frac12 (S_+^1 S_-^2 + S_-^1 S_+^2).
\]
Nous avons donc l'hamiltonien :
\[
H_H =
J\left[
\frac12(S_+^1 S_-^2 + S_-^1 S_+^2)
S_z^1 S_z^2
\right].
\]
\[
J\hbar^2
\begin{pmatrix}
\frac14 & 0 & 0 & 0\\
0 & -\frac14 & \frac12 & 0\\
0 & \frac12 & -\frac14 & 0\\
0 & 0 & 0 & \frac14\\
\end{pmatrix}.
\]

\subsection*{(c) Énergies et états propres de $H_H$}
Dans la base
$\{\ket{\uparrow\uparrow},\ket{\uparrow\downarrow},\ket{\downarrow\uparrow},\ket{\downarrow\downarrow}\}$,
l'hamiltonien Heisenberg s'écrit (comme vu en b)),
\[
H_H = J\hbar^2
\begin{pmatrix}
\frac14 & 0 & 0 & 0\\[4pt]
0 & -\frac14 & \frac12 & 0\\[4pt]
0 & \frac12 & -\frac14 & 0\\[4pt]
0 & 0 & 0 & \frac14
\end{pmatrix}.
\]
Les états \(\ket{\uparrow\uparrow}\) et \(\ket{\downarrow\downarrow}\) sont immédiatement propres avec valeur propre \(J\hbar^2/4\).
Pour le sous-espace engendré par \(\{\ket{\uparrow\downarrow},\ket{\downarrow\uparrow}\}\) on diagonalise le bloc
\[
J\hbar^2\begin{pmatrix}-\tfrac14 & \tfrac12\\[4pt]\tfrac12 & -\tfrac14\end{pmatrix}.
\]
En résolvant \((-\tfrac14-\lambda)x + \tfrac12 y = 0\) et \(\tfrac12 x + (-\tfrac14-\lambda)y = 0\), on trouve les valeurs propres
\(\lambda_1=\tfrac14\) et \(\lambda_2=-\tfrac34\).
Les vecteurs propres respectifs sont :
\[
\ket{\phi_1}=\frac{1}{\sqrt2}\big(\ket{\uparrow\downarrow}+\ket{\downarrow\uparrow}\big),\qquad
\ket{\phi_2}=\frac{1}{\sqrt2}\big(\ket{\uparrow\downarrow}-\ket{\downarrow\uparrow}\big).
\]
Ainsi, les états propres et leurs énergies associées sont :
\[
\begin{aligned}
\ket{\phi_0} &= \ket{\uparrow\uparrow}, & E_0 &= \dfrac{J\hbar^2}{4},\\[4pt]
\ket{\phi_1} &= \dfrac{1}{\sqrt2}\big(\ket{\uparrow\downarrow}+\ket{\downarrow\uparrow}\big), & E_1 &= \dfrac{J\hbar^2}{4},\\[4pt]
\ket{\phi_2} &= \ket{\downarrow\downarrow}, & E_2 &= \dfrac{J\hbar^2}{4},\\[6pt]
\ket{\phi_3} &= \dfrac{1}{\sqrt2}\big(\ket{\uparrow\downarrow}-\ket{\downarrow\uparrow}\big), & E_3 &= -\dfrac{3J\hbar^2}{4}.
\end{aligned}
\]

\subsection*{(d) Application d'un champ magnétique}

On exprime le champ magnétique t.q. :
$H_B = -\gamma B_0\big(S_1^z+S_2^z\big)$.\\

Calculons le commutateur :
\[
[H_H,H_B] = -\gamma B_0\sum_{\mu=x,y,z}\Big([S_\mu^1 S_\mu^1,S_z^1+[S_\mu^1 S_\mu^2,S_z^2]\Big).
\]
En utilisant la propriété \([S_a^i,S_b^j]=i\hbar\,\varepsilon_{abc}\delta_{ij}S_c^i\) on obtient :
\[
[S_\mu^1 S_\mu^2,S_z^1] = i\hbar\,\varepsilon_{\mu z c}\,S_c^1 S_\mu^2,\qquad
[S_\mu^1 S_\mu^2,S_z^2] = i\hbar\,\varepsilon_{\mu z c}\,S_\mu^1 S_2^c.
\]
La somme sur les indices \(\mu,c\) de \(\varepsilon_{\mu z c}(S_1^c S_2^\mu + S_1^\mu S_2^c)\) s'annule par antisymétrie, donc :
\([H_H,S_1^z+S_2^z]=0\) et donc \([H_H,H_B]=0\).

\paragraph{Conclusion}
On conclue que comme \(H_H\) commute avec \(H_B\), ils sont simultanément diagonalisables : les états propres de \(H_H\) peuvent être choisis comme états propres de \(H_B\). \(H_B\) est donc diagonal dans la base \(\{\ket{\uparrow\uparrow},\ket{\uparrow\downarrow},\ket{\downarrow\uparrow},\ket{\downarrow\downarrow}\}\) :
\[
S_z^1+S_z^2=\hbar\,\mathrm{diag}(1,0,0,-1),
\quad
H_B=-\gamma B_0\hbar\,\mathrm{diag}(1,0,0,-1).
\]

Les états propres et énergies de \(H'=H_H+H_B\) sont donc :
\[
\begin{aligned}
\ket{\phi_0}=\ket{\uparrow\uparrow},\quad &E_{\phi_0}'=\tfrac{J\hbar^2}{4}-\gamma B_0\hbar,\\[4pt]
\ket{\phi_1}=\tfrac{1}{\sqrt2}(\ket{\uparrow\downarrow}+\ket{\downarrow\uparrow}),\quad &E_{\phi_1}'=\tfrac{J\hbar^2}{4},\\[4pt]
\ket{\phi_2}=\ket{\downarrow\downarrow},\quad &E_{\phi_2}'=\tfrac{J\hbar^2}{4}+\gamma B_0\hbar,\\[6pt]
\ket{\phi_3}=\tfrac{1}{\sqrt2}(\ket{\uparrow\downarrow}-\ket{\downarrow\uparrow}),\quad &E_{\phi_3}'=-\tfrac{3J\hbar^2}{4}.
\end{aligned}
\]

Le champ magnétique retire la dégénérescence qu'il y avait entre $\ket{\phi_0}$, $\ket{\phi_1}$ et $\ket{\phi_2}$.\\
$\ket{\phi3}$, pour sa part, reste inchangé.


\subsection*{(e) Terme DM : construction complète}
En simplifiant les vecteurs $S_1$ et $S_2$, on obtient que l'hamiltonien peut être défini t.q :
\[
H_{DM} =D(S_x^1 S_y^2 - S_y^1 S_x^2)
\]

On peut exprimer $S_x$ et $S_y$ par leur définition trouvée en a) :
\[
S_x = \frac{S_+ + S_-}{2},\qquad
S_y = \frac{S_+ - S_-}{2i}.
\]

Développons les termes avec leur définition :
\[
S_x^1 S_y^2 =
\frac{1}{4i}(S_+^1 + S_-^1)(S_+^2 - S_-^2),
\]
\[
S_y^1 S_x^2 =
\frac{1}{4i}(S_+^1 - S_-^1)(S_+^2 + S_-^2).
\]
\[
= \frac{1}{2i}(S_-^1 S_+^2 - S_+^1 S_-^2).
\]

On met le tout sous forme de matrices :
\[
\hbar^2
\begin{pmatrix}
0 & 0 & 0 & 0\\
0 & 0 & -1 & 0\\
0 & 1 & 0 & 0\\
0 & 0 & 0 & 0\\
\end{pmatrix}.
\]

\[
\frac{D\hbar^2}{2}
\begin{pmatrix}
0 & 0 & 0 & 0\\
0 & 0 & i & 0\\
0 & -i & 0 & 0\\
0 & 0 & 0 & 0\\
\end{pmatrix}.
\]

Les valeurs propres du bloc central sont 
$\lambda_{\pm} = \pm \frac{D\hbar^2}{2}$.

\[
\boxed{
E_\pm = \pm \frac{D\hbar^2}{2},\qquad
\ket{\phi_\pm} = \frac{1}{\sqrt2}
(\ket{\uparrow\downarrow} \pm i\ket{\downarrow\uparrow}).
}
\]

Les états $\ket{\uparrow\uparrow}$ et $\ket{\downarrow\downarrow}$ sont d'énergie 0.
\end{question}

\newpage

\begin{question}{Produit tensoriel de deux spins $\frac{1}{2}$}

\subsection*{(a)Valeurs propres et états propres de H dans $\sigma_z$}
On rappelle que l'hamiltonien est donné par $H=\hbar g\,\sigma_{y}^1\otimes\sigma_{x}^2$.\\

Nous avons nos matrices de pauli données par :
\[
\sigma_x=\begin{pmatrix}0&1\\1&0\end{pmatrix},\quad
\sigma_y=\begin{pmatrix}0&-i\\i&0\end{pmatrix},\quad
\sigma_z=\begin{pmatrix}1&0\\0&-1\end{pmatrix}.
\]

On obtient en faisant le produit tensoriel que la matrice de $H$ est t.q. :
\[
\sigma_y\otimes\sigma_x =
\begin{pmatrix}
0 & 0 & 0 & -i\\[4pt]
0 & 0 & -i & 0\\[4pt]
0 & i & 0 & 0\\[4pt]
i & 0 & 0 & 0
\end{pmatrix},\qquad
H=\hbar g\,\sigma_y\otimes\sigma_x.
\]


Puisque $(\sigma_y\otimes\sigma_x)^2=I_4$, les valeurs propres de $\sigma_y\otimes\sigma_x$ sont $\pm1$. Les valeurs propres de $H$ sont donc
\[
E=\pm\hbar g.
\]

Les états propres de $\sigma_x$ et $\sigma_y$ sont donnés par:
\[
\ket{+}_x=\frac{1}{\sqrt2}\begin{pmatrix}1\\1\end{pmatrix},\;
\ket{-}_x=\frac{1}{\sqrt2}\begin{pmatrix}1\\-1\end{pmatrix},\;
\ket{+}_y=\frac{1}{\sqrt2}\begin{pmatrix}1\\ i\end{pmatrix},\;
\ket{-}_y=\frac{1}{\sqrt2}\begin{pmatrix}1\\ -i\end{pmatrix}.
\]
Les possibles combinaisons avec le produit tensoriel sont alors vecteurs propres ; en base computationnelle :
\[
\begin{aligned}
\ket{+_y,+_x}
&= \frac{1}{2}\begin{pmatrix}1\\[2pt]1\\[2pt]i\\[2pt]i\end{pmatrix},
\qquad
\ket{+_y,-_x}
= \frac{1}{2}\begin{pmatrix}1\\[2pt]-1\\[2pt]i\\[2pt]-i\end{pmatrix},\\[6pt]
\ket{-_y,+_x}
&= \frac{1}{2}\begin{pmatrix}1\\[2pt]1\\[2pt]-i\\[2pt]-i\end{pmatrix},
\qquad
\ket{-_y,-_x}
= \frac{1}{2}\begin{pmatrix}1\\[2pt]-1\\[2pt]-i\\[2pt]i\end{pmatrix}.
\end{aligned}
\]

Valeurs propres associées :
\[
\begin{aligned}
&H\ket{+_y,+_x}=+\hbar g\,\ket{+_y,+_x},\quad
H\ket{+_y,-_x}=-\hbar g\,\ket{+_y,-_x},\\[4pt]
&H\ket{-_y,+_x}=-\hbar g\,\ket{-_y,+_x},\quad
H\ket{-_y,-_x}=+\hbar g\,\ket{-_y,-_x}.
\end{aligned}
\]

\paragraph{(b) Mesure d'énergie pour } $\ket{\psi}=\dfrac{1}{\sqrt2}(\ket{\uparrow\uparrow}+\ket{\downarrow\downarrow})$.\\
On réécrit l'état en vecteur de coordonnées t.q. :
\[
\ket{\psi}=\frac{1}{\sqrt2}\begin{pmatrix}1\\[2pt]0\\[2pt]0\\[2pt]1\end{pmatrix}.
\]
On projette $\ket{\psi}$ sur la base $\{\ket{+_y,+_x}, \ket{-_y,-_x}, \ket{-_y,+_x}, \ket{+_y,-_x}\}$ : 

\[
\begin{aligned}
a_{+,+} &= \braket{+_y,+_x|\psi}
= \frac{1}{2\sqrt2}(1-i),\\[6pt]
a_{+,-} &= \braket{+_y,-_x|\psi}
= \frac{1}{2\sqrt2}(1+i),\\[6pt]
a_{-,+} &= \braket{-_y,+_x|\psi}
= \frac{1}{2\sqrt2}(1+i),\\[6pt]
a_{-,-} &= \braket{-_y,-_x|\psi}
= \frac{1}{2\sqrt2}(1-i).
\end{aligned}
\]
Calcul des probabilités :
\[
|a|^2 = \frac{1}{8}\,|1\pm i|^2 = \frac{1}{8}\cdot 2 = \frac14 \quad\text{pour chaque paire }(+_{xy},-_{xy}).
\]
Deux états propres donnent $E=+\hbar g$ et deux donnent $E=-\hbar g$, donc
\[
\boxed{P(E=+\hbar g)=\tfrac12,\qquad P(E=-\hbar g)=\tfrac12.}
\]

\paragraph{(c) Mesure en fonction des probabilités}
Si l'on mesure d'abord $S_{z}$ et qu'on obtient $+\tfrac{\hbar}{2}$ (projection sur $\ket{\uparrow}_1$), l'état projeté est :
\[
(\ket{\uparrow}\bra{\uparrow}\otimes I)\ket{\psi} = \frac{1}{\sqrt2}\ket{\uparrow\uparrow}.
\]
L'état normalisé du second spin est $\ket{\uparrow}_2$. La probabilité que la mesure suivante $S_{x}$ donne $+\tfrac{\hbar}{2}$ est
\[
P(S_{x}=\braket{+\tfrac{\hbar}{2}\mid S_{z}}=+\tfrac{\hbar}{2})
= |\braket{+_x|\uparrow}|^2 = \left(\frac{1}{\sqrt2}\right)^2 = \frac12.
\]
\end{question}

\newpage

\begin{question}{Oscillateur harmonique chargé dans un champ électrique}
Petit rappel :
\[
H_0=\frac{P^2}{2m}+\frac12 m\omega^2 X^2 = \hbar\omega\Big(a^\dagger a + \tfrac12\Big).
\]
La perturbation est :
\[
W(t) = -q\varepsilon(t)X.
\]
Le nouvel hamiltonien est donné par :
\[
H = H_0 + W
\]
\subsection*{(a) Expression de (H(t)) en ($a,a^\dagger$)}
Remplaçons (X) par son expression en ($a,a^\dagger$) :
\[
X = \sqrt{\frac{\hbar}{2m\omega}}(a+a^\dagger).
\]
Donc
\[
W(t)=-q\varepsilon(t)\sqrt{\frac{\hbar}{2m\omega}}(a+a^\dagger) = -\hbar\lambda(t)(a+a^\dagger),
\]
où $\lambda(t)=\dfrac{q\varepsilon(t)}{\sqrt{2\hbar m\omega}}$.

\vspace{0.5cm}

Ainsi le nouvel hamiltonien est :
\[
\boxed{H(t)=\hbar\omega\Big(a^\dagger a + \tfrac12\Big) - \hbar\lambda(t)\big(a+a^\dagger\big)}
\]


\subsection*{(b) Calcul des commutateurs ($[a,H]$) et ($[a^\dagger,H]$)}

On utilise $[a,a^\dagger]=1$ et $[a,a]=[a^\dagger,a^\dagger]=0$.\\

Développons $H(t)=\hbar\omega a^\dagger a + \tfrac12\hbar\omega - \hbar\lambda(t)(a+a^\dagger)$.

\paragraph{Calcul de $[a,\hbar\omega a^\dagger a]$.}
Utilisons la relation :
\[
[a,a^\dagger a] = [a,a^\dagger]a + a^\dagger[a,a] = 1a + a^\dagger\cdot 0 = a.
\]
Donc
\[
[a,\hbar\omega a^\dagger a]=\hbar\omega a.
\]

\paragraph{Calcul de $[a,\tfrac12\hbar\omega]$.}
Pusique $h\omega$ est une constate scalaire, le commutateur est nul donc :
\[
[a,\tfrac12\hbar\omega]=0.
\]

\paragraph{Calcul de $[a,-\hbar\lambda(t)(a+a^\dagger)]$.}
\[
[a,-\hbar\lambda(t)(a+a^\dagger)]
= -\hbar\lambda(t)\big([a,a] + [a,a^\dagger]\big)
= -\hbar\lambda(t)\big(0 + 1\big)
= -\hbar\lambda(t).
\]

\[
\boxed{[a,H] = \hbar\omega a - \hbar\lambda(t).}
\]

\medskip

On fait la même chose pour $a^\dagger$.
\paragraph{Calcul de $[a^\dagger,\hbar\omega a^\dagger a]$.}
\[
[a^\dagger,a^\dagger a] = a^\dagger[a^\dagger,a] + [a^\dagger,a^\dagger]a = a^\dagger(-1) + 0 = -a^\dagger,
\]
donc,
\[
[a^\dagger,\hbar\omega a^\dagger a] = \hbar\omega(-a^\dagger) = -\hbar\omega a^\dagger.
\]
\vspace{0.5cm}
Même chose pour le calcul de $[a^\dagger,\tfrac12\hbar\omega]$, le commutateur donne 0.

\paragraph{Calcul de $[a^\dagger,-\hbar\lambda(t)(a+a^\dagger)]$.}
\[
[a^\dagger, -\hbar\lambda(t)(a+a^\dagger)]
= -\hbar\lambda(t)\big([a^\dagger,a] + [a^\dagger,a^\dagger]\big)
= -\hbar\lambda(t)\big(-1 + 0\big)
= +\hbar\lambda(t).
\]

\[
\boxed{[a^\dagger,H] = -\hbar\omega a^\dagger + \hbar\lambda(t).}
\]

\subsection*{(c) Preuve que $\alpha(t)$ satisfait $\dfrac{d\alpha (t)}{dt} = -i\omega\alpha (t) + i\lambda (t)$}

En mécanique quantique, pour un opérateur O sans dépendance explicite en t , l'équation d'Ehrenfest est :
\[
\frac{d}{dt}\langle O\rangle = \frac{1}{i\hbar}\langle [O,H]\rangle.
\]
Appliquons ceci à $O=a$ :
\[
\dot\alpha(t) = \frac{1}{i\hbar}\langle [a,H]\rangle.
\]
En remplaçant $[a,H]$ trouvé précédemment :
\[
\dot\alpha(t) = \frac{1}{i\hbar}\langle \hbar\omega a - \hbar\lambda(t)\rangle
= \frac{1}{i}\big(\omega \alpha(t) - \lambda(t)\big).
\]
On trouve :
\[
\boxed{\displaystyle \dot\alpha(t) = -i\omega\alpha(t) + i\lambda(t).}
\]

\subsection*{(d) Intégration de l'équation trouvée en c)}

Réécrivons l'équation :
\[
\dot\alpha(t) + i\omega \alpha(t) = i\lambda(t).
\]
\vspace{0.25cm}
On pose que $\mu(t)=e^{i\omega t}$ puisque $\frac{d}{dt}e^{i\omega t} = i\omega e^{i\omega t}$. Multiplions l'équation par $\mu(t)$ :
\[
e^{i\omega t}\dot\alpha(t) + i\omega e^{i\omega t}\alpha(t) = i e^{i\omega t}\lambda(t).
\]
Le côté gauche est la dérivée totale :
\[
\frac{d}{dt}\big(e^{i\omega t}\alpha(t)\big) = i e^{i\omega t}\lambda(t).
\]
Intégrons de $t_0$ à $t$ :
\[
e^{i\omega t}\alpha(t) - e^{i\omega t_0}\alpha_0 = i\int_{t_0}^{t} e^{i\omega t'}\lambda(t')dt'.
\]
Isolons $\alpha(t)$ :
\[
\boxed{\displaystyle
\alpha(t) = \alpha_0 e^{-i\omega(t-t_0)} + i\int_{t_0}^{t} e^{-i\omega(t-t')}\lambda(t')dt'.
}
\]

\subsection*{(e) Expressions de $\langle X\rangle$ et $\langle P\rangle$ en fonction de $\alpha(t)$}

Rappel des opérateurs :
\[
X = \sqrt{\frac{\hbar}{2m\omega}}(a+a^\dagger),\qquad
P = i\sqrt{\frac{\hbar m\omega}{2}}(a^\dagger-a).
\]

En prenant la valeur moyenne de chaque opérateur, et en posant $\alpha(t)=\langle a\rangle$, on obtient
\[
\boxed{\displaystyle
\langle X\rangle(t) = \sqrt{\frac{\hbar}{2m\omega}}\big(\alpha(t)+\alpha(t)\big),
}
\]
\[
\boxed{\displaystyle
\langle P\rangle(t) = i\sqrt{\frac{\hbar m\omega}{2}}\big(\alpha(t)-\alpha(t)\big).
}
\]
Pour un état cohérent initial $\ket{\psi(t)}$ évoluant sous l'Hamiltonien libre $H_0$, les moyennes $\langle X\rangle(t)$ et $\langle P\rangle(t)$ valent des fonctions cosinus/sinus de fréquence $\omega$. En l'absence de champ électrique, la trajectoire moyenne $\langle X\rangle(t)$ obéit à l'équation différentielle d'un oscillateur harmonique libre : le mouvement est donc harmonique.
\end{question}
\end{document}